\section{平面関数と臨界条件}
標数零で考える。$D=x_1^3+x_2^3+x_3^3$とし、巡回する$i,j,k$に対して
\begin{equation}\label{mh:family}
 F_{s,P,Q,T}=D/3+u^3+v^3+s x_1x_2x_3+
 \sum_i(P_i u+Q_i v)x_jx_k+\sum_i T_i x_iuv
\end{equation}
を取る。$u=\ell=\sum\lambda_i x_i$, $v=m=\sum\mu_i x_i$は平面の局所座標である。
原点に合同な三変数三次式$f$には、一意な$N\equiv I$と$\sigma$が存在して
$f(Nx)=D/3+\sigma x_1x_2x_3$となる。
実際、九つの行列方向の微分は$x_i^2x_j$の九単項式であり、残る$xyz$方向と合わせて全三次式の基底をなす。
形式逆関数定理を用いれば、全Artin基底に両立する一意な規格化を得る。
この規格化を制限式へ適用した係数を$\sigma(s,P,Q,T,\lambda,\mu)$と書く。
$\sigma(s,P,Q,T,0,0)=s$であり、全変数の通常次数で
\begin{equation}\label{mh:jet}
 \sigma=s+\sum_i(P_i\lambda_i+Q_i\mu_i)
 +6(\lambda_1\lambda_2\lambda_3+\mu_1\mu_2\mu_3)
 +\sum_{\{i,j,k\}=\{1,2,3\}}T_k\lambda_i\mu_j+O(4).
\end{equation}
二次の規格化は非対角行列であり、その$sxyz$への作用に$xyz$項はないので、この三次jetは規格化後にも成立する。

\begin{proposition}[内在的臨界条件]\label{mh:intrinsic}
滑らかな平面三次曲線$E=V(f)\subset W\subset X=V(F)$に対し、平面移動による曲線の変形が零となる条件は
\[
 F_u|_W,F_v|_W\in (\partial f)_2.
\]
これは極化写像$V\to\Sym^2 W^*$の階数が三であること、
および$0\to T_E\to T_X|_E\to N_{E/X}\simeq\mathcal O_E(1)^2\to0$の分裂と同値である。
このとき$V=W\oplus\ker DF_W$という標準補空間を取り、
\[
 F=f(x)+\sum_i x_iq_i(u,v)+g(u,v)
\]
と書ける。特に平方線形混合項を含まない。
\end{proposition}
\begin{proof}
三つの偏微分は二次式の正則列であり、Jacobian環のHilbert関数は$(1,3,3,1)$である。
その完全対合$R_1\otimes R_2\to R_3$から、任意の一次式$a,b$に対し
$aF_u+bF_v$がJacobianイデアルに入る条件は、表示の二つの二次類の消滅と等しい。
これは平面切断のKodaira--Spencer写像の消滅である。
同じ写像は法束列の連結写像であり、$T_E\simeq\mathcal O_E$とSerre双対を用いると、
その消滅は拡張類の消滅に等しい。三つの$x$偏微分は独立なので極化の階数は少なくとも三であり、
残る二偏微分がその像に入ることが階数三と等しい。
核は$W$と交わらず、核方向を$u,v$に選べば$x^2u,x^2v$項が消える。
\end{proof}

以後、臨界点での三次式を
\begin{equation}\label{mh:criticalfamily}
 F_Q=D/3+u^3+v^3+\sum_i x_i(a_i u^2+d_iuv+c_i v^2),\qquad
 Q_i=\begin{pmatrix}2a_i&d_i\\d_i&2c_i\end{pmatrix}
\end{equation}
とする。平面変数を交互順序$n=(\lambda_1,\mu_1,\lambda_2,\mu_2,\lambda_3,\mu_3)$に並べると、
$\sigma$のHessianおよび三次項は
\begin{equation}\label{mh:hessian}
 H=\begin{pmatrix}0&Q_3&Q_2\\Q_3&0&Q_1\\Q_2&Q_1&0\end{pmatrix},\qquad
 \Gamma(n)=6(\lambda_1\lambda_2\lambda_3+\mu_1\mu_2\mu_3).
\end{equation}
入力の平面次数が二以上なので、Hesse規格化の非線形効果は四次以上であり、これは正確な二、三次項である。

\section{六次Hessian因子}
\begin{theorem}\label{mh:delta}
九係数空間で$\Delta=\det H$と置くと
\[
 \Delta=4\prod_i(4a_ic_i-d_i^2)-16\det(a,d,c)^2.
\]
$\mathcal D=V(\Delta)\subset\mathbb A^9$は絶対既約な正規有理八次元超曲面で、
特異集合の次元は六である。原点での完備局所環も正規整域である。
階数五で$\ker H=\langle(n_1,n_2,n_3)\rangle$, $n_i\in K^2$と書けば、
$\mathcal D$が特異である条件は非零な$n_i$が高々一つであることである。
\end{theorem}
\begin{proof}
対称二行行列$A,B,C$について、直接Schur補元を取ると
\[
 \det\begin{pmatrix}0&C&B\\C&0&A\\B&A&0\end{pmatrix}
 =4\det A\det B\det C-
 \det\begin{pmatrix}A_{11}&A_{12}&A_{22}\\B_{11}&B_{12}&B_{22}\\C_{11}&C_{12}&C_{22}\end{pmatrix}^{\!2}.
\]
$C$可逆の稠密開集合で合同変換により$C=I$とし、さらに$A$を対角化すると、
右辺は$\det(AB+BA)$の計算に帰着する。多項式恒等式なので全体に延長する。
$Q_i$を代入すると表示式を得る。

$A=\left(\begin{smallmatrix}x&y\\y&z\end{smallmatrix}\right)$,
$B=\left(\begin{smallmatrix}u&v\\v&w\end{smallmatrix}\right)$,
$C=\left(\begin{smallmatrix}r&s\\s&t\end{smallmatrix}\right)$とする。
上式を$R=\CC[x,y,z,u,v,w]$上の$r,s,t$に関する二次式と見る。
$r^2,t^2$係数はそれぞれ$-(yw-zv)^2,-(xv-yu)^2$であり、二つの一次原始既約式は非同伴なので二次式は原始的である。
$A=I,B=\operatorname{diag}(1,2)$で二次式は$8rt-9s^2$となり、その行列式は$144$である。
従って$\operatorname{Frac}R$の代数閉包上でも階数三で、一次式二つの積にはならない。
Gaussの補題から絶対既約性を得る。

$A$可逆の開集合では$\operatorname{rank}H=4+\operatorname{rank}(CA^{-1}B+BA^{-1}C)$。
$A=I$に帰着し、$B=\left(\begin{smallmatrix}x&y\\y&z\end{smallmatrix}\right)$に対する
$C\mapsto BC+CB$の行列は
\[
 \begin{pmatrix}2x&2y&0\\y&x+z&y\\0&2y&2z\end{pmatrix},\qquad
 \det=4(xz-y^2)(x+z).
\]
その階数が一以下なら二次小行列式から$B=0$となる。
よって階数三、二、零の層で$H$の階数四以下の集合の次元はいずれも高々六である。
三行列とも特異な補集合も次元高々六である。
階数五では$\operatorname{adj}H=\kappa kk^t$なので、$Q_i$方向の微分は
$2\kappa n_j^t(\delta Q_i)n_k$である。
任意の対称行列$S$について$e^tSf=0$となるのは$e=0$または$f=0$の場合だけであり、特異条件を得る。
例えば$n_1\ne0,n_2=n_3=0$なら$Q_2n_1=Q_3n_1=0$であり、
核直線の一次元、$Q_1$の三次元、$Q_2,Q_3$それぞれの一次元を足して六次元以下である。
一方$Q_1=0$の六次元空間は階数四以下なので全体が特異である。
従って特異集合の次元は正確に六であり、整超曲面の$(S_2)$と余次元二以上の特異集合から$(R_1)$を得て正規性が従う。
$\Delta$は斉次既約なので完備環の随伴次数環も$\CC[Q]/(\Delta)$という整域である。
分離的フィルトレーションから完備環も整域であり、excellent正規局所環の完備化の正規性を用いる。

有理性には核incidenceを用いる。$[k]\in\PP^5$に対する
$L_k:\Sym^2(K^2)^{\oplus3}\to K^6$, $Q\mapsto H(Q)k$の全射開集合上で、
incidenceは階数三のベクトル束である。
\[
 k=(1,0,0,1,1,1)^t,\quad
 Q_1=\begin{pmatrix}1&1\\1&-1\end{pmatrix},\quad
 Q_2=\begin{pmatrix}-1&1\\1&2\end{pmatrix},\quad
 Q_3=\begin{pmatrix}-2&0\\0&-3\end{pmatrix}
\]
では$L_k$の列$1,2,3,4,5,8$の行列式は$2$、$H$の第一行列を除く小行列式は$-12$である。
階数五の上では核直線が一意なので、この有理八次元束は$\mathcal D$と双有理である。
\end{proof}

\section{核上三次の四成分}
この節の行列成分は$Q_i=\left(\begin{smallmatrix}A_i&B_i\\B_i&C_i\end{smallmatrix}\right)$と記す。
$\mathcal D_5=\{\Delta=0,\operatorname{rank}H=5\}$上で
$\Gamma|_{\ker H}=0$の零スキームを$\mathcal C_5$、そのスキーム論的閉包を$\overline{\mathcal C}_5$とする。
\begin{theorem}\label{mh:four}
$\overline{\mathcal C}_5$は被約純七次元で、ちょうど四つの有理的既約成分
$M,E_1,E_2,E_3$を持つ。$T=\det(A,B,C)$と置くと
\[
 E_i=V(\det Q_i,T),\qquad
 \deg\PP(M)=27,\quad\deg\PP(E_i)=6.
\]
$M$は三ブロック全て非零な核を持つ部分の閉包である。
随伴行列の第$j$列を$v_j$とすると全閉包の有限イデアル表示は
\[
 (\Delta,\Gamma(v_1),\ldots,\Gamma(v_6)):I_5(H)^\infty.
\]
\end{theorem}
\begin{proof}
核incidenceの階数をまず全て数える。非零ブロックが三つなら$\operatorname{rank}L_k=6$、
二つで独立なら五、二つで比例なら四、一つなら四であり、対応する射影基底の次元は$5,3,2,1$である。
例えば独立な二ブロックを$e_1,e_2$、第三を$(r,s)$にした六方程式では、六次小行列式に$-2r,-2s$が現れる。
三つとも比例する場合は$n_i=t_ie_1$とすれば二組の三次行列式が$2t_1t_2t_3$である。
ただしこのとき全ての$Q_i e_1=0$で$H$の階数は三以下である。
二つで独立かつ第三零なら$Q_3=0$となり$H$の階数は四以下である。
残る場合の核条件は$Q_3n=0,t_1Q_2n+t_2Q_1n=0$、
または$Q_2n_1=Q_3n_1=0$であり、表示の階数が従う。

$X=V(\lambda_1\lambda_2\lambda_3+\mu_1\mu_2\mu_3)\subset\PP^5$は原始一次式として既約であり、
$\lambda_1\lambda_2\ne0$上で有理的である。特異集合は$\lambda_i$座標点と$\mu_j$座標点を結ぶ九本の直線である。
従って正規で、三ブロック非零の開集合$U$上では滑らかである。
$X\cap U$上の階数三束は七次元の有理的主成分を与える。
核の一ブロックが零で他の二つが比例する各場合は、二次元の有理的基底上の階数五束を与え、七次元成分となる。
一ブロックのみ非零の部分は六次元以下である。
$\mathcal D_5$はCohen--Macaulayで、$\Gamma$は零でない切断なのでその零スキームは純七次元である。
従って以上で成分を尽くす。

例えば$Q_3$階数一なら合同変換で$Q_3=\operatorname{diag}(0,e)$とでき、
$T=e(A_1B_2-B_1A_2)$である。これから$V(\det Q_3,T)$の$Q_3\ne0$部分は既約で一般点被約、
$Q_3=0$部分は六次元である。二、三次式の完全交差なので埋込み成分はなく、全体が整で第三境界成分に一致する。
他も同様である。
主成分では滑らかな$X\cap U$により一般点の重複度は一である。
境界では
\[
 Q_1=\begin{pmatrix}1&1\\1&2\end{pmatrix},\quad
 Q_2=\begin{pmatrix}-1&-1\\-1&3\end{pmatrix},\quad
 Q_3=\begin{pmatrix}0&0\\0&1\end{pmatrix},\quad k=(1,0,1,0,0,0)
\]
を用いる。第一行列を除く小行列$J$は$\det J=-8$、$(J^{-1})_{\lambda_3,\lambda_3}=-1/8$。
$\delta A_1=1$は$T_Q\mathcal D$に属し、正規化した核の微分は$\delta\lambda_3=1/8$、
従って$d\Gamma=3/4\ne0$である。Cohen--Macaulay性から全体が被約である。
階数五チャート上では随伴行列の列が核の生成元を張るため、飽和イデアルの式も従う。

次数は$\PP^8_Q\times\PP^5_k$上の六本の双次数$(1,1)$方程式$Hk=0$で計算する。
上の階数表からincidenceは純七次元の整完全交差であり、その類は$(A+B)^6$である。
$\Gamma=0$との交差を$Q$側に押すと$3B(A+B)^6$の$B^5$係数は$45A^2$である。
階数四以下に写る部分は像の射影次元が五以下なので寄与しない。
三境界成分はそれぞれ次数六、各成分の写像次数・重複度は一なので、主成分は$45-18=27$となる。
\end{proof}

\begin{proposition}[一ブロック核の非孤立極限]\label{mh:nonisolated}
階数五で核が$(n,0,0)$なら、平面関数は形式右同値で$\sigma=q_5(y)$となる。
従って$h^2=3\sigma$のファイバーは五次元通常二重点と形式直線の積であり、
$T^1=K[[z]]$である。
\end{proposition}
\begin{proof}
$Q_2n=Q_3n=0$である。$u=z\lambda x_1,v=z\mu x_1$への制限は
$c(z)x_1^3+(x_2^3+x_3^3)/3$、$c(0)=1/3$となる。
対角規格化でHesse係数は恒等的に零であり、二つの正規偏微分も$x_1^2$の倍数なので全平面微分が零である。
高速五方向の停留解は一意なのでこの曲線に一致する。形式Morse分裂により残余関数は厳密に零である。
\end{proof}
